\documentclass{standalone}
\usepackage{circuitikz}
\usetikzlibrary{calc}
\definecolor{circuitnotered}{HTML}{E5534B}
\definecolor{circuitnotemark}{HTML}{FE00FE}
\begin{document}
\begin{circuitikz}[american, line width=0.8pt]
\ctikzset{bipoles/length=1.2cm}
\foreach \x in {0,2,4} {\foreach \y in {0,-2,-4,-6,-8} {\fill[gray, opacity=0.35] (\x,\y) circle (1.1pt);}}
\node[gray, font=\scriptsize] at (0,0.84) {1};
\node[gray, font=\scriptsize] at (2,0.84) {2};
\node[gray, font=\scriptsize] at (4,0.84) {3};
\node[gray, font=\scriptsize, anchor=east] at (-0.84,0) {a};
\node[gray, font=\scriptsize, anchor=east] at (-0.84,-2) {b};
\node[gray, font=\scriptsize, anchor=east] at (-0.84,-4) {c};
\node[gray, font=\scriptsize, anchor=east] at (-0.84,-6) {d};
\node[gray, font=\scriptsize, anchor=east] at (-0.84,-8) {e};
\coordinate (a1) at (0,0);
\coordinate (a2) at (2,0);
\draw (a1) to[R, l_=$R_{1}$, a^=$10\,\mathrm{k}\Omega$] (a2); % line 3
\draw[circuitnotered] (2,-2) circle (0.9); % line 5
\path (2,-2.593) rectangle (6.001,-1.407); % line 6
\node[anchor=west, inner sep=0, circuitnotemark, font=\large] at (2,-2) {X}; % line 6
\draw[circuitnotered] (2,-4) circle (0.9); % line 7
\path (-0.477,-4.593) rectangle (4.477,-3.407); % line 8
\node[anchor=west, inner sep=0, circuitnotemark, font=\large] at (2,-4) {X}; % line 8
\draw[circuitnotered] (2,-6) circle (0.9); % line 9
\path (-2.254,-6.593) rectangle (2,-5.407); % line 10
\node[anchor=west, inner sep=0, circuitnotemark, font=\large] at (2,-6) {X}; % line 10
\path (2,-8.593) rectangle (5.937,-7.407); % line 11
\node[anchor=west, inner sep=0, circuitnotemark, font=\large\bfseries] at (2,-8) {X}; % line 11
\path (5.407,-1.905) rectangle (6.593,0); % line 12
\node[anchor=west, inner sep=0, circuitnotemark, font=\large] at (6,0) {X}; % line 12
\path (7.407,-8) rectangle (8.593,-5.396); % line 13
\node[anchor=west, inner sep=0, circuitnotemark, font=\large] at (8,-8) {X}; % line 13
\path (12,-9.356) rectangle (22.279,0.593); % line 14
\node[anchor=west, inner sep=0, circuitnotemark, font=\large] at (12,0) {X}; % line 14
\node[anchor=west, inner sep=0, circuitnotemark, font=\large] at (12,-0.487) {X}; % line 14
\node[anchor=west, inner sep=0, circuitnotemark, font=\large] at (12,-0.974) {X}; % line 14
\node[anchor=west, inner sep=0, circuitnotemark, font=\large] at (12,-1.46) {X}; % line 14
\node[anchor=west, inner sep=0, circuitnotemark, font=\large] at (12,-1.947) {X}; % line 14
\node[anchor=west, inner sep=0, circuitnotemark, font=\large] at (12,-2.434) {X}; % line 14
\node[anchor=west, inner sep=0, circuitnotemark, font=\large] at (12,-2.921) {X}; % line 14
\node[anchor=west, inner sep=0, circuitnotemark, font=\large] at (12,-3.408) {X}; % line 14
\node[anchor=west, inner sep=0, circuitnotemark, font=\large] at (12,-3.895) {X}; % line 14
\node[anchor=west, inner sep=0, circuitnotemark, font=\large] at (12,-4.381) {X}; % line 14
\node[anchor=west, inner sep=0, circuitnotemark, font=\large] at (12,-4.868) {X}; % line 14
\node[anchor=west, inner sep=0, circuitnotemark, font=\large] at (12,-5.355) {X}; % line 14
\node[anchor=west, inner sep=0, circuitnotemark, font=\large] at (12,-5.842) {X}; % line 14
\node[anchor=west, inner sep=0, circuitnotemark, font=\large] at (12,-6.329) {X}; % line 14
\node[anchor=west, inner sep=0, circuitnotemark, font=\large] at (12,-6.816) {X}; % line 14
\node[anchor=west, inner sep=0, circuitnotemark, font=\large] at (12,-7.302) {X}; % line 14
\node[anchor=west, inner sep=0, circuitnotemark, font=\large] at (12,-7.789) {X}; % line 14
\node[anchor=west, inner sep=0, circuitnotemark, font=\large] at (12,-8.276) {X}; % line 14
\node[anchor=west, inner sep=0, circuitnotemark, font=\large] at (12,-8.763) {X}; % line 14
\node[anchor=south west, inner sep=2pt, circuitnotemark, font=\LARGE\bfseries] (circuittitle) at (current bounding box.north west) {X};
\path (circuittitle.south west) rectangle ++(7.498,0.853);
\end{circuitikz}
\end{document}
